%================================================================
\chapter{Python Virtual Environments and Workflow Automation}
\label{ch:python-environments}
%================================================================

Before we write any scripts, we need to understand \emph{what} those scripts
will manage: Python virtual environments.

\textbf{The Problem: Package Conflicts}.
Imagine you have two projects on your computer. Project~A uses version~1.5
of a library called \texttt{numpy}; Project~B requires version~2.0. If you
install both versions globally, they will conflict. The last one installed
wins, and one of your projects could break.

\textbf{The Solution: Virtual Environments}.
A \textbf{virtual environment} (or \emph{venv}) is an isolated copy of
Python with its own set of installed packages, completely separate from
every other environment on your machine. Think of it as a self-contained
bubble around your project.

\begin{itemize}
  \item Each project gets its own environment
  \item Packages installed in one environment cannot interfere with another
  \item You can recreate an environment exactly from a list of packages
  \item Deleting an environment is safe and clean
\end{itemize}

We will cover two strategies for managing environments: The traditional 
approach via \texttt{pip} and an emerging trend called \texttt{uv}.

%---------------------------------------------------------
\section{Traditional Package Management: \texttt{pip}}
\label{sec:venvs}
%---------------------------------------------------------
\subsubsection{Windows (CMD or PowerShell)}

\begin{lstlisting}[style=windowsStyle, caption={Creating a venv on Windows},escapeinside={(*@}{@*)}]
:: Navigate to your project folder
cd C:\Projects\(*@\ProjectName@*)

:: Create the environment (Python must be installed globally)
python -m venv C:\penv\(*@\ProjectName@*)
\end{lstlisting}

\subsubsection{Linux / WSL}

\begin{lstlisting}[style=linuxStyle, caption={Creating a venv on Linux/WSL},escapeinside={(*@}{@*)}]
# Navigate to your project folder
cd ~/(*@\ProjectName@*)

# Create the environment
python3 -m venv ~/penv/(*@\ProjectName@*)
\end{lstlisting}

In both cases, \texttt{C:\textbackslash penv\textbackslash \ProjectName{}} or
\texttt{\textasciitilde/penv/\ProjectName{}} is where the environment will be stored.
Keeping all environments in a single folder (\texttt{penv} stands for
``Python environments'') makes them easy to find and manage.

\subsection{Activating a Virtual Environment}

Activating an environment modifies your current shell session so that
\texttt{python} and \texttt{pip} refer to the isolated copies inside the
venv, not the global ones.

\subsubsection{Windows}

\begin{lstlisting}[style=windowsStyle, caption={Activating a venv on Windows},escapeinside={(*@}{@*)}]
C:\penv\(*@\ProjectName@*)\Scripts\activate.bat
\end{lstlisting}

After activation, your prompt will change to show the environment name:
\begin{lstlisting}[style=windowsStyle,escapeinside={(*@}{@*)}]
((*@\ProjectName@*)) C:\Projects\(*@\ProjectName@*)>
\end{lstlisting}

\subsubsection{Linux / WSL}

\begin{lstlisting}[style=linuxStyle, caption={Activating a venv on Linux/WSL},escapeinside={(*@}{@*)}]
source ~/penv/(*@\ProjectName@*)/bin/activate
\end{lstlisting}

After activation:
\begin{lstlisting}[style=linuxStyle,escapeinside={(*@}{@*)}]
((*@\ProjectName@*)) drg@machine:~/(*@\ProjectName@*)$
\end{lstlisting}

\begin{warningbox}
Note the \texttt{source} keyword in the Linux command. This is critical and
is explained in detail in Section~\ref{sec:source-vs-execute}.
\end{warningbox}

\subsection{Installing Packages}

Once the environment is active, use \texttt{pip} to install packages.
The syntax is identical on both platforms:

\begin{lstlisting}[style=inlineStyle, caption={Installing packages with pip},escapeinside={(*@}{@*)}]
pip install numpy pandas matplotlib openai langchain
\end{lstlisting}

To save a list of all installed packages (so you or a colleague can
recreate the environment later):

\begin{lstlisting}[style=inlineStyle,escapeinside={(*@}{@*)}]
pip freeze > requirements.txt
\end{lstlisting}

To recreate an environment from such a list:

\begin{lstlisting}[style=inlineStyle,escapeinside={(*@}{@*)}]
pip install -r requirements.txt
\end{lstlisting}

\subsection{Deactivating}

To leave the environment and return to the global shell (identical on
both platforms):

\begin{lstlisting}[style=inlineStyle,escapeinside={(*@}{@*)}]
deactivate
\end{lstlisting}

\subsection{Installing the \ProjectName{} Required Libraries}
\label{sec:project-libraries}

With the \ProjectName{} environment active, install the required libraries for
the data challenge:

\begin{verbatim}
pip install numpy pandas matplotlib scipy jupyter
\end{verbatim}

These five libraries are sufficient for all analysis in the data
challenge. Additional libraries required for the LLM chapter (including
\texttt{openai}, \texttt{anthropic}, and \texttt{langchain}) are
installed separately in Chapter~\ref{ch:ollama-infrastructure}.

\vs
\noindent
To verify the installation:

\begin{verbatim}
python -c "import numpy, pandas, matplotlib, scipy, jupyter; print('OK')"
\end{verbatim}

\vs
\noindent
\textbf{You know your core environment is fully functional when:}
\begin{itemize}
    \item \texttt{python --version} returns a version number.
    \item \texttt{code .} opens Visual Studio Code from a terminal.
    \item \texttt{git --version} returns a version number.
    \item The five required libraries import without errors (see above).
\end{itemize}
The LLM agent programs (\texttt{agentGroq.py}, \texttt{agentGPT.py},
\texttt{agentClaude.py}) are located at
\texttt{resources/unit\_7/FOO/} in the \ProjectName{} repository.
Running them requires API keys configured in
Chapter~\ref{ch:ollama-infrastructure}.
\vs

% ============================================================
\section{Modern Package Management: \texttt{uv}}
\label{sec:uv}
% ============================================================

\subsection{Why \texttt{uv} Matters Today}

The tools introduced in the previous section (\texttt{python -m venv}
for environment creation and \texttt{pip} for package
installation) are the established, stable foundation of Python package
management. They are maintained by the Python Packaging Authority
(PyPA) under the umbrella of the Python Software Foundation (PSF), a
non-profit organisation with a long-term governance structure and no
commercial dependency. Every Python installation includes \texttt{pip}
by default, and knowing how to use it remains an essential skill
regardless of what other tools you adopt.

A newer tool called \texttt{uv} has emerged and is rapidly gaining
adoption in professional environments. Understanding both tools is
important precisely because the landscape is in transition.

\begin{infobox}
\textbf{A note on governance.} \texttt{pip} is maintained by the
Python Packaging Authority under the Python Software Foundation, a
non-profit with a community governance model and no commercial
dependency. \texttt{uv} is an open-source tool developed and
maintained by Astral, a venture-capital funded commercial company.
Astral has been an excellent steward of the tool and has contributed
meaningfully to the broader Python packaging ecosystem. Nevertheless,
its long-term trajectory depends on business outcomes in a way that
foundation-backed projects do not. Until the ecosystem stabilises
around a clear community standard, knowing both \texttt{pip} and
\texttt{uv} is the prudent professional position.
\end{infobox}

\subsection{What \texttt{uv} Offers}

\texttt{uv} is a drop-in replacement for \texttt{pip} and
\texttt{python -m venv}, rewritten in Rust for performance. Its
principal advantages over the traditional toolchain are:

\begin{itemize}
  \item \textbf{Speed.} Package resolution and installation are
        10--100 times faster than \texttt{pip} due to a Rust
        implementation and aggressive parallel downloading.
  \item \textbf{A central cache.} Downloaded packages are stored once
        in a global cache and reused across environments. On the same
        filesystem, \texttt{uv} uses hardlinks so that multiple
        environments sharing a package consume no additional disk space
        beyond the first copy.
  \item \textbf{Lockfiles.} \texttt{uv} produces a \texttt{uv.lock}
        file that records the complete, resolved dependency graph with
        cryptographic hashes. This is a significant improvement over
        \texttt{pip freeze > requirements.txt}, which records only
        installed package names and versions without distinguishing
        direct dependencies from transitive ones, and without hash
        verification.
  \item \textbf{\texttt{pyproject.toml} integration.} \texttt{uv}
        works natively with \texttt{pyproject.toml}, the modern
        standard for declaring project dependencies, replacing the
        older \texttt{requirements.txt} workflow.
  \item \textbf{Python version management.} \texttt{uv} can download
        and manage multiple Python versions independently, removing the
        need for separate tools such as \texttt{pyenv}.
\end{itemize}

\subsection{Installing \texttt{uv}}

\subsubsection{Windows (PowerShell)}

\begin{lstlisting}[style=windowsStyle, caption={Installing \texttt{uv} on Windows},escapeinside={(*@}{@*)}]
powershell -ExecutionPolicy ByPass -c "irm https://astral.sh/uv/install.ps1 | iex"
\end{lstlisting}

\texttt{uv} installs itself to
\texttt{\%USERPROFILE\%\textbackslash.local\textbackslash bin} and
adds itself to your user \texttt{PATH} automatically. Close and reopen
your terminal, then verify:

\begin{lstlisting}[style=windowsStyle,escapeinside={(*@}{@*)}]
uv --version
\end{lstlisting}

\subsubsection{Linux / WSL}

\begin{lstlisting}[style=linuxStyle, caption={Installing \texttt{uv} on Linux/WSL},escapeinside={(*@}{@*)}]
curl -LsSf https://astral.sh/uv/install.sh | sh
\end{lstlisting}

\texttt{uv} installs to \texttt{\textasciitilde/.local/bin}. Reload
your shell configuration and verify:

\begin{lstlisting}[style=linuxStyle,escapeinside={(*@}{@*)}]
source ~/.bashrc
uv --version
\end{lstlisting}

\subsection{Creating and Managing Environments with \texttt{uv}}

The environment structure produced by \texttt{uv venv} is identical to
a standard \texttt{python -m venv} environment. The same
\texttt{Scripts\textbackslash activate.bat} (Windows) and
\texttt{bin/activate} (Linux) scripts are present and activation is
unchanged.

\subsubsection{Creating an Environment}

\begin{lstlisting}[style=windowsStyle, caption={Creating a uv-managed environment (Windows)},escapeinside={(*@}{@*)}]
uv venv C:\uvenv\(*@\ProjectName@*) --python 3.12
\end{lstlisting}

\begin{lstlisting}[style=linuxStyle, caption={Creating a uv-managed environment (Linux/WSL)},escapeinside={(*@}{@*)}]
uv venv ~/uvenv/(*@\ProjectName@*) --python 3.12
\end{lstlisting}

The \texttt{--python} flag specifies the Python version. If that
version is not already installed, \texttt{uv} will download and
install it automatically.

\subsubsection{Activating an Environment}

Activation is identical to a standard venv:

\begin{lstlisting}[style=windowsStyle,escapeinside={(*@}{@*)}]
C:\uvenv\(*@\ProjectName@*)\Scripts\activate.bat
\end{lstlisting}

\begin{lstlisting}[style=linuxStyle,escapeinside={(*@}{@*)}]
source ~/uvenv/(*@\ProjectName@*)/bin/activate
\end{lstlisting}

\subsubsection{Installing Packages}

Within an activated environment, \texttt{uv pip install} replaces
\texttt{pip install}:

\begin{lstlisting}[style=inlineStyle, caption={Installing packages with uv},escapeinside={(*@}{@*)}]
uv pip install numpy pandas matplotlib
\end{lstlisting}

\subsubsection{Declaring Dependencies with \texttt{pyproject.toml}}

For projects that will be shared or reproduced, \texttt{uv} works best
with a \texttt{pyproject.toml} file in the project root that declares
dependencies explicitly:

\begin{lstlisting}[style=inlineStyle, caption={Minimal \texttt{pyproject.toml}},escapeinside={(*@}{@*)}]
[project]
name = "my-project"
version = "0.1.0"
requires-python = ">=3.12"
dependencies = [
    "numpy>=1.26",
    "pandas>=2.2",
    "matplotlib>=3.9",
]
\end{lstlisting}

With this file present, a single command installs all dependencies and
generates the lockfile:

\begin{lstlisting}[style=inlineStyle,escapeinside={(*@}{@*)}]
uv sync
\end{lstlisting}

\texttt{uv sync} reads \texttt{pyproject.toml}, resolves the full
dependency graph, installs all packages, and writes \texttt{uv.lock}.
Running \texttt{uv sync} again on another machine, or after a
\texttt{git pull} that added new dependencies, reproduces the
environment exactly.

\begin{infobox}
Commit \texttt{pyproject.toml} and \texttt{uv.lock} to your
repository. Do not commit the environment folder itself. Add it to
\texttt{.gitignore}.
\end{infobox}

\newpage
\subsection{\texttt{requirements.txt} vs.\ \texttt{uv.lock}}

\begin{table}[h]
\centering
\caption{Comparison of \texttt{requirements.txt} and \texttt{uv.lock}}
\label{tab:req-vs-lock}
\begin{tabular}{@{}p{0.28\textwidth}p{0.30\textwidth}p{0.30\textwidth}@{}}
\toprule
\textbf{Property} & \textbf{\texttt{requirements.txt}} & \textbf{\texttt{uv.lock}} \\
\midrule
Generated by
  & \texttt{pip freeze}
  & \texttt{uv sync} / \texttt{uv lock} \\[4pt]
Records
  & Installed packages only
  & Full dependency graph \\[4pt]
Distinguishes direct vs.\ transitive deps
  & No
  & Yes \\[4pt]
Hash verification
  & No
  & Yes \\[4pt]
Cross-platform reproducibility
  & Partial
  & Yes \\[4pt]
Should be committed to git
  & Yes
  & Yes \\
\bottomrule
\end{tabular}
\end{table}

% ============================================================
\section{Level 1: Script Files}
\label{sec:level1}
% ============================================================

Now that we understand virtual environments, we can write scripts to
activate them automatically. The goal is to create two scripts:

\begin{enumerate}
  \item \textbf{\texttt{ActivateEnv}}: A reusable script that accepts
        two parameters (environment name and project directory), activates
        the environment, and opens VS Code.
  \item \textbf{\texttt{\ProjectName}}: A project-specific one-liner that calls
        \texttt{ActivateEnv} with the correct parameters for the \ProjectName{}
        project.
\end{enumerate}

\subsection{The Windows Solution: Batch Files}

\subsubsection{ActivateEnv.bat}

The Windows version, \texttt{ActivateEnv.bat}, is stored in
\texttt{C:\textbackslash penv} and placed on the system \texttt{PATH}.
It accepts two parameters: \texttt{\%1} (environment name) and
\texttt{\%2} (project directory).

\begin{lstlisting}[style=windowsStyle, basicstyle=\footnotesize\ttfamily, caption={\texttt{ActivateEnv.bat}: Reusable environment launcher (Windows)},escapeinside={(*@}{@*)}]
@echo off
setlocal

:: Read parameters
set "ENV_NAME=%~1"
set "BASE_DIR=%~2"
echo Starting activation of environment '\%ENV_NAME%'...

:: Validate input
if "%ENV_NAME%"=="" (
    echo [ERROR] Environment name not specified.
    goto :eof
)
if "%BASE_DIR%"=="" (
    echo [ERROR] Base directory not specified.
    goto :eof
)

:: Check if base directory exists
if not exist "%BASE_DIR%" (
    echo [ERROR] Directory not found: %BASE_DIR%
    goto :eof
)

:: Change to working directory
cd /d "%BASE_DIR%"

:: Define activation script path
set "VENV_ROOT=C:\penv"
set "VENV_ACTIVATE=%VENV_ROOT%\%ENV_NAME%\Scripts\activate.bat"

:: Activate the environment
if exist "%VENV_ACTIVATE%" (
    echo Activating...
    call "%VENV_ACTIVATE%"
    echo Activation complete.
) else (
    echo [ERROR] Not found: %VENV_ACTIVATE%
    goto :eof
)

:: Open VS Code in the project directory
echo Launching VS Code...
code .

endlocal
\end{lstlisting}


Key Windows-specific elements:
\begin{itemize}
  \item \texttt{@echo off} suppresses the echoing of each command to the
        screen.
  \item \texttt{setlocal} ensures that variable changes are local to this
        script.
  \item \texttt{\%\textasciitilde1} strips surrounding quotes from the first
        argument.
  \item \texttt{call} is used to invoke another batch file and return
        control afterward.
  \item \texttt{goto :eof} jumps to the end of the file (a clean exit).
\end{itemize}

\subsubsection{\ProjectName{}.bat}

This file lives in the \ProjectName{} project directory. It simply calls
\texttt{ActivateEnv.bat} with the correct arguments:

\begin{lstlisting}[style=windowsStyle, caption={\texttt{\ProjectName{}.bat}: Project launcher (Windows)},escapeinside={(*@}{@*)}]
ActivateEnv (*@\ProjectName@*) "C:\Projects\(*@\ProjectName@*)"
\end{lstlisting}

Because \texttt{C:\textbackslash penv} is on the system \texttt{PATH},
Windows finds \texttt{ActivateEnv.bat} automatically. You can run
\texttt{\ProjectName{}.bat} by double-clicking it or typing \texttt{\ProjectName{}} in any
CMD window.

\subsection{The Linux Solution: Shell Scripts}
\label{subsec:linux-scripts}

\subsubsection{Understanding Shell Script Syntax}

A Linux shell script is a plain text file containing Bash commands. Every
script should begin with a \textbf{shebang line}, a special first line that
tells the system which interpreter to use:

\begin{lstlisting}[style=linuxStyle,escapeinside={(*@}{@*)}]
#!/bin/bash
\end{lstlisting}

Parameters passed to a script are accessed as \texttt{\$1}, \texttt{\$2},
\texttt{\$3}, etc. The script's own name is \texttt{\$0}.

\subsubsection{ActivateEnv.sh}

\begin{lstlisting}[style=linuxStyle, basicstyle=\footnotesize\ttfamily, caption={\texttt{ActivateEnv.sh}: Reusable environment launcher (Linux/WSL)},escapeinside={(*@}{@*)}]
#!/bin/bash
# ActivateEnv.sh
# Usage: source ~/penv/ActivateEnv.sh <ENV_NAME> <PROJECT_DIR>
# Note: Must be *sourced*, not executed. See explanation below.

ENV_NAME="$1"
BASE_DIR="$2"

echo "Starting activation of environment '${ENV_NAME}'..."

# Validate input
if [ -z "$ENV_NAME" ]; then
    echo "[ERROR] Environment name not specified."
    return 1
fi

if [ -z "$BASE_DIR" ]; then
    echo "[ERROR] Base directory not specified."
    return 1
fi

# Check if the project directory exists
if [ ! -d "$BASE_DIR" ]; then
    echo "[ERROR] Directory not found: ${BASE_DIR}"
    return 1
fi

# Change to the project directory
cd "$BASE_DIR" || return 1

# Build the path to the activation script
VENV_ROOT="$HOME/penv"
VENV_ACTIVATE="${VENV_ROOT}/${ENV_NAME}/bin/activate"

# Activate the environment
if [ -f "$VENV_ACTIVATE" ]; then
    echo "Activating..."
    # shellcheck source=/dev/null
    source "$VENV_ACTIVATE"
    echo "Activation complete."
else
    echo "[ERROR] Not found: ${VENV_ACTIVATE}"
    return 1
fi

# Open VS Code in the project directory (WSL-aware)
echo "Launching VS Code..."
code .
\end{lstlisting}

Notice several differences from the Windows version:
\begin{itemize}
  \item \texttt{-z} tests whether a string is empty (zero length).
  \item \texttt{-d} tests whether a path is a directory.
  \item \texttt{-f} tests whether a path is a regular file.
  \item \texttt{\$HOME} is the Linux equivalent of
        \texttt{\%USERPROFILE\%}: it expands to your home directory.
  \item \texttt{return 1} exits the script with an error code (we use
        \texttt{return} rather than \texttt{exit} for a reason explained in
        the next section).
\end{itemize}

\subsubsection{\ProjectName{}.sh}

\begin{lstlisting}[style=linuxStyle, caption={\texttt{\ProjectName.sh}: Project launcher (Linux/WSL)},escapeinside={(*@}{@*)}]
#!/bin/bash
# (*@\ProjectName@*).sh
# Usage: source ~/(*@\ProjectName@*)/(*@\ProjectName@*).sh
# Launches the (*@\ProjectName@*) project environment.

source ~/penv/ActivateEnv.sh (*@\ProjectName@*) "$HOME/(*@\ProjectName@*)"
\end{lstlisting}

\subsection{Making a Script Executable}
\label{sec:chmod}

On Linux, creating a file does not automatically make it executable. You
must grant execute permission using the \texttt{chmod} command:

\begin{lstlisting}[style=linuxStyle,escapeinside={(*@}{@*)}]
chmod +x ~/penv/ActivateEnv.sh
chmod +x ~/(*@\ProjectName@*)/(*@\ProjectName@*).sh
\end{lstlisting}

\texttt{chmod} stands for \emph{change mode}. The \texttt{+x} flag adds
execute (\texttt{x}) permission. Linux permissions are organized around
three roles (owner, group, and others), but for personal scripts,
\texttt{chmod +x} is almost always all you need.

\subsection{Adding Your Scripts to PATH}

On Windows, you added \texttt{C:\textbackslash penv} to the system
\texttt{PATH} so that \texttt{ActivateEnv.bat} could be found from
anywhere. On Linux, you do the same by editing your shell configuration
file.

\begin{lstlisting}[style=linuxStyle, caption={Adding \textasciitilde/penv to PATH in \textasciitilde/.bashrc},escapeinside={(*@}{@*)}]
# Open your bashrc file in a text editor
nano ~/.bashrc

# Add this line at the bottom:
export PATH="$HOME/penv:$PATH"

# Save and reload the configuration
source ~/.bashrc
\end{lstlisting}

After this, the shell can find \texttt{ActivateEnv.sh} from any directory.

% ------------------------------------------------------------
\subsection{The \texttt{uv} Approach: \texttt{ActivateUVEnv}}
\label{subsec:activateuvenv}
% ------------------------------------------------------------

The following scripts mirror \texttt{ActivateEnv} but target
\texttt{uv}-managed environments stored in
\texttt{C:\textbackslash uvenv} (Windows) or
\texttt{\textasciitilde/uvenv} (Linux). Store all \texttt{uv}
environments in \texttt{uvenv} rather than \texttt{penv} to keep the
two toolchains visually distinct on disk.

Two behaviours distinguish \texttt{ActivateUVEnv} from its
\texttt{pip}-based counterpart:

\begin{enumerate}
  \item \textbf{Automatic creation.} If the named environment does not
        yet exist, \texttt{uv venv} creates it before activation. There
        is no separate setup step for new projects.
  \item \textbf{Automatic synchronisation.} If a
        \texttt{pyproject.toml} is present in the project directory,
        \texttt{uv sync} runs after activation. Dependencies install
        themselves after every \texttt{git pull} that adds new
        requirements, so no manual \texttt{pip install} is needed.
\end{enumerate}

\subsubsection{ActivateUVEnv.bat (Windows)}

Store this file in \texttt{C:\textbackslash uvenv} and add that folder
to the system \texttt{PATH}.

\begin{lstlisting}[style=windowsStyle, basicstyle=\footnotesize\ttfamily, caption={\texttt{ActivateUVEnv.bat}: uv environment launcher (Windows)},escapeinside={(*@}{@*)}]
@echo off
setlocal

set "ENV_NAME=%~1"
set "BASE_DIR=%~2"
set "VENV_ROOT=C:\uvenv"
set "VENV_PATH=%VENV_ROOT%\%ENV_NAME%"
set "VENV_ACTIVATE=%VENV_PATH%\Scripts\activate.bat"

:: Validate inputs
if "%ENV_NAME%"=="" (
    echo [ERROR] Usage: ActivateUVEnv ^<env-name^> ^<project-dir^>
    goto :eof
)
if "%BASE_DIR%"=="" (
    echo [ERROR] Usage: ActivateUVEnv ^<env-name^> ^<project-dir^>
    goto :eof
)
if not exist "%BASE_DIR%" (
    echo [ERROR] Project directory not found: %BASE_DIR%
    goto :eof
)

:: Create environment if it does not exist yet
if not exist "%VENV_ACTIVATE%" (
    echo [INFO] Environment '%ENV_NAME%' not found. Creating with uv...
    uv venv "%VENV_PATH%" --python 3.12
    if errorlevel 1 (
        echo [ERROR] Failed to create environment.
        goto :eof
    )
)

:: Activate
cd /d "%BASE_DIR%"
call "%VENV_ACTIVATE%"

:: Sync dependencies if pyproject.toml is present
if exist "pyproject.toml" (
    echo [INFO] pyproject.toml found. Syncing dependencies...
    uv sync
)

echo.
echo [OK] Environment '%ENV_NAME%' active.
echo Launching VS Code...
code .

endlocal
\end{lstlisting}

Usage is identical in form to \texttt{ActivateEnv.bat}:

\begin{lstlisting}[style=windowsStyle,escapeinside={(*@}{@*)}]
ActivateUVEnv lunar "C:\Dropbox\Research\Lunar\python"
ActivateUVEnv forecasting "C:\Projects\forecasting"
\end{lstlisting}

\subsubsection{ActivateUVEnv: Linux / WSL Shell Function}

Add the following function to \texttt{\textasciitilde/.bashrc}. Like
its \texttt{pip}-based counterpart, it runs in the current shell so
that activation and directory changes persist in your session.

\begin{lstlisting}[style=linuxStyle, basicstyle=\footnotesize\ttfamily, caption={\texttt{ActivateUVEnv()} shell function for \textasciitilde/.bashrc},escapeinside={(*@}{@*)}]
# -------------------------------------------------------
# ActivateUVEnv: Activate a uv-managed venv and open VS Code
# Usage: ActivateUVEnv <ENV_NAME> <PROJECT_DIR>
# -------------------------------------------------------
ActivateUVEnv() {
    local ENV_NAME="$1"
    local BASE_DIR="$2"
    local VENV_ROOT="$HOME/uvenv"
    local VENV_PATH="${VENV_ROOT}/${ENV_NAME}"
    local VENV_ACTIVATE="${VENV_PATH}/bin/activate"

    if [ -z "$ENV_NAME" ]; then
        echo "[ERROR] Usage: ActivateUVEnv <env-name> <project-dir>"
        return 1
    fi

    if [ -z "$BASE_DIR" ]; then
        echo "[ERROR] Usage: ActivateUVEnv <env-name> <project-dir>"
        return 1
    fi

    if [ ! -d "$BASE_DIR" ]; then
        echo "[ERROR] Project directory not found: ${BASE_DIR}"
        return 1
    fi

    # Create environment if it does not exist yet
    if [ ! -f "$VENV_ACTIVATE" ]; then
        echo "[INFO] Environment '${ENV_NAME}' not found. Creating with uv..."
        uv venv "$VENV_PATH" --python 3.12 || return 1
    fi

    # Activate
    cd "$BASE_DIR" || return 1
    # shellcheck source=/dev/null
    source "$VENV_ACTIVATE"

    # Sync dependencies if pyproject.toml is present
    if [ -f "pyproject.toml" ]; then
        echo "[INFO] pyproject.toml found. Syncing dependencies..."
        uv sync
    fi

    echo "[OK] Environment '${ENV_NAME}' active."
    echo "Launching VS Code..."
    code .
}
\end{lstlisting}

Reload \texttt{.bashrc} after saving:

\begin{lstlisting}[style=linuxStyle,escapeinside={(*@}{@*)}]
source ~/.bashrc
\end{lstlisting}

Then call it from any directory:

\begin{lstlisting}[style=linuxStyle,escapeinside={(*@}{@*)}]
ActivateUVEnv lunar "$HOME/Research/Lunar/python"
\end{lstlisting}

% ============================================================
\section{The Critical Difference: Executing vs.\ Sourcing}
\label{sec:source-vs-execute}
% ============================================================

This section explains one of the most important, and most
misunderstood, concepts in Linux shell scripting, especially for
developers coming from Windows.

\subsection{Child Shells}

Every time you run a script with \texttt{./script.sh} or just
\texttt{script.sh}, Linux launches a \textbf{child shell}: a new,
separate shell process that inherits a copy of your environment, runs
the script, and then \emph{terminates}. Any changes the script makes to
the environment (activating a venv, changing directory, setting a
variable) exist only in that child shell. When the child terminates,
those changes disappear. Your original terminal is completely unaffected.

This is why running \texttt{./ActivateEnv.sh} would seem to do
nothing useful: the venv gets activated inside the child, then the child
exits, and your terminal is still pointing at the global Python.

\subsection{Sourcing: Running in the Current Shell}

The \texttt{source} command (or its shorthand, a single dot \texttt{.})
runs a script \emph{in the current shell}: no child is created. Every
command executes as if you had typed it directly at the prompt. This means:

\begin{itemize}
  \item \texttt{source ~/penv/\ProjectName{}/bin/activate} activates the venv
        in \emph{your} terminal.
  \item \texttt{cd} changes \emph{your} working directory.
  \item Variables set by the script persist in \emph{your} session.
\end{itemize}

\begin{lstlisting}[style=linuxStyle,escapeinside={(*@}{@*)}]
# These two are equivalent:
source ActivateEnv.sh (*@\ProjectName@*) ~/(*@\ProjectName@*)
. ActivateEnv.sh (*@\ProjectName@*) ~/(*@\ProjectName@*)

# This does NOT work for venv activation:
./ActivateEnv.sh (*@\ProjectName@*) ~/(*@\ProjectName@*)    # child shell; changes are lost
\end{lstlisting}

\subsection{The Windows Parallel}

Windows does not have this distinction in quite the same way. The
\texttt{call} command in a batch file runs another \texttt{.bat} in the
same CMD session and returns, which is why the Windows version works
as expected. Linux's default behavior of spawning a child process is
actually the \emph{safer} design, as it prevents scripts from accidentally
modifying your environment, but it requires you to opt in to the
shared-session behavior with \texttt{source}.

\begin{notebox}
This is also why \texttt{ActivateEnv.sh} uses \texttt{return} instead of
\texttt{exit} for error handling. When a script is sourced, \texttt{exit}
would close your entire terminal. \texttt{return} exits only the script,
leaving your session intact.
\end{notebox}

% ============================================================
\section{Level 2: Shell Functions, The Elegant Linux Approach}
\label{sec:level2}
% ============================================================

Sourcing scripts is powerful, but it requires typing \texttt{source}
or \texttt{.} every time. It also means your launcher script must live
somewhere on the \texttt{PATH} and the user must remember the exact
invocation. Linux offers a more elegant solution: \textbf{shell functions}.

\subsection{What Is a Shell Function?}

A shell function is like a script, but it is defined directly inside
your shell configuration file (\texttt{\textasciitilde/.bashrc}).
Because it is part of the shell itself (not a separate file), it
automatically runs in the current session without requiring
\texttt{source}. You call it just by typing its name.

\subsection{Defining the ActivateEnv Function}

Add the following to the bottom of \texttt{\textasciitilde/.bashrc}:

\begin{lstlisting}[style=linuxStyle, basicstyle=\footnotesize\ttfamily, caption={Shell function version of \texttt{ActivateEnv} in \textasciitilde/.bashrc},escapeinside={(*@}{@*)}]
# -------------------------------------------------------
# ActivateEnv: Activate a Python venv and open VS Code
# Usage: ActivateEnv <ENV_NAME> <PROJECT_DIR>
# -------------------------------------------------------
ActivateEnv() {
    local ENV_NAME="$1"
    local BASE_DIR="$2"

    echo "Starting activation of environment '${ENV_NAME}'..."

    # Validate input
    if [ -z "$ENV_NAME" ]; then
        echo "[ERROR] Environment name not specified."
        return 1
    fi

    if [ -z "$BASE_DIR" ]; then
        echo "[ERROR] Base directory not specified."
        return 1
    fi

    if [ ! -d "$BASE_DIR" ]; then
        echo "[ERROR] Directory not found: ${BASE_DIR}"
        return 1
    fi

    local VENV_ACTIVATE="$HOME/penv/${ENV_NAME}/bin/activate"

    if [ ! -f "$VENV_ACTIVATE" ]; then
        echo "[ERROR] Virtual environment not found: ${VENV_ACTIVATE}"
        return 1
    fi

    # Change directory, activate, and open VS Code
    cd "$BASE_DIR" || return 1
    source "$VENV_ACTIVATE"
    echo "Environment '${ENV_NAME}' activated."
    code .
}
\end{lstlisting}

The keyword \texttt{local} restricts variables to the function's scope,
so they do not pollute your shell session.

\subsection{Defining the \ProjectName{} Function}

Now add a one-line function that calls \texttt{ActivateEnv} for the
\ProjectName{} project:

\begin{lstlisting}[style=linuxStyle, caption={Project-specific function for \ProjectName{} in \textasciitilde/.bashrc},escapeinside={(*@}{@*)}]
# -------------------------------------------------------
# (*@\ProjectName@*): Launch the (*@\ProjectName@*) project environment
# Usage: (*@\ProjectName@*)
# -------------------------------------------------------
(*@\ProjectName@*)() {
    ActivateEnv "(*@\ProjectName@*)" "$HOME/(*@\ProjectName@*)"
}
\end{lstlisting}

\subsection{Activating the New Functions}

After editing \texttt{\textasciitilde/.bashrc}, reload it:

\begin{lstlisting}[style=linuxStyle,escapeinside={(*@}{@*)}]
source ~/.bashrc
\end{lstlisting}

Now, from \emph{any directory}, simply type:

\begin{lstlisting}[style=linuxStyle,escapeinside={(*@}{@*)}]
(*@\ProjectName@*)
\end{lstlisting}

The shell will change to your \ProjectName{} project directory, activate the
virtual environment, and open VS Code, all in one word, with no
\texttt{source} prefix required.

\subsection{Why Windows Has No Clean Equivalent}

Windows batch files always run as child processes; there is no
\texttt{source} command in CMD. PowerShell comes closest: you can define
functions in your PowerShell profile (\texttt{\$PROFILE}), and they
behave similarly to Bash functions. However, PowerShell's profile is
disabled by default for security reasons and requires changing the
execution policy. For Windows users who want this one-command experience
natively, configuring PowerShell profiles is the path forward, but it
involves more administrative overhead than editing
\texttt{\textasciitilde/.bashrc}. This is one of the practical reasons
many developers on Windows adopt WSL.

% ============================================================
\section{Putting It All Together}
\label{sec:together}
% ============================================================

Figure~\ref{fig:workflow} summarizes the complete workflow on each platform.

\begin{table}[h]
\centering
\caption{Complete workflow comparison: Windows vs.\ Linux/WSL}
\label{fig:workflow}
\begin{tabular}{@{}p{0.15\textwidth}p{0.38\textwidth}p{0.38\textwidth}@{}}
\toprule
\textbf{Step} & \textbf{Windows} & \textbf{Linux / WSL} \\
\midrule
Create venv
  & \texttt{python -m venv C:\textbackslash penv\textbackslash \ProjectName}
  & \texttt{python3 -m venv \textasciitilde/penv/\ProjectName} \\[4pt]
Reusable launcher
  & \texttt{ActivateEnv.bat} in \texttt{C:\textbackslash penv}
  & Function \texttt{ActivateEnv()} in \texttt{\textasciitilde/.bashrc} \\[4pt]
Project launcher
  & \texttt{\ProjectName{}.bat} in project folder
  & Function \texttt{\ProjectName{}()} in \texttt{\textasciitilde/.bashrc} \\[4pt]
Make available globally
  & Add \texttt{C:\textbackslash penv} to system PATH
  & Functions in \texttt{.bashrc} are always available \\[4pt]
Launch project
  & Type \texttt{\ProjectName{}} in any CMD window
  & Type \texttt{\ProjectName{}} in any terminal \\[4pt]
Result
  & venv active, VS Code opens in project folder
  & venv active, VS Code opens in project folder (WSL mode) \\
\bottomrule
\end{tabular}
\end{table}

\subsection{The Full WSL Experience}

When you type \texttt{\ProjectName{}} in your WSL terminal:

\begin{enumerate}
  \item Bash looks up \texttt{\ProjectName{}} in its list of known functions.
  \item \texttt{\ProjectName{}()} calls \texttt{ActivateEnv "\ProjectName{}" "\$HOME/\ProjectName{}"}.
  \item \texttt{ActivateEnv()} changes to the project directory.
  \item The virtual environment is activated in your current shell.
  \item \texttt{code .} launches VS Code with the WSL extension,
        connecting your Windows VS Code interface to your Linux backend.
  \item You are now working in a fully configured development environment
        with one word.
\end{enumerate}

% ============================================================
\section{Quick Reference Cheat Sheet}
\label{sec:cheatsheet}
% ============================================================

\begin{table}[h]
\centering
\caption{Shell scripting and virtual environment cheat sheet}
\label{tab:cheatsheet}
\begin{tabular}{@{}p{0.35\textwidth}p{0.28\textwidth}p{0.28\textwidth}@{}}
\toprule
\textbf{Task} & \textbf{Windows (CMD)} & \textbf{Linux / WSL (Bash)} \\
\midrule
Create virtual environment
  & \texttt{python -m venv C:\textbackslash penv\textbackslash NAME}
  & \texttt{python3 -m venv \textasciitilde/penv/NAME} \\[4pt]
Activate environment
  & \texttt{C:\textbackslash penv\textbackslash NAME\textbackslash Scripts\textbackslash activate}
  & \texttt{source \textasciitilde/penv/NAME/bin/activate} \\[4pt]
Deactivate environment
  & \texttt{deactivate}
  & \texttt{deactivate} \\[4pt]
Install a package
  & \texttt{pip install pkg}
  & \texttt{pip install pkg} \\[4pt]
Save package list
  & \texttt{pip freeze > req.txt}
  & \texttt{pip freeze > req.txt} \\[4pt]
Restore from list
  & \texttt{pip install -r req.txt}
  & \texttt{pip install -r req.txt} \\[4pt]
Make script executable
  & (not required; extension handles it)
  & \texttt{chmod +x script.sh} \\[4pt]
Run a script
  & \texttt{script.bat}
  & \texttt{./script.sh} \\[4pt]
Source a script
  & \texttt{call script.bat}
  & \texttt{source script.sh} \\[4pt]
Edit shell config
  & (no equivalent in CMD)
  & \texttt{nano \textasciitilde/.bashrc} \\[4pt]
Reload shell config
  & (restart CMD)
  & \texttt{source \textasciitilde/.bashrc} \\[4pt]
Open VS Code (WSL)
  & \texttt{code .}
  & \texttt{code .} \\[4pt]
Script parameter~1
  & \texttt{\%\textasciitilde1}
  & \texttt{\$1} \\[4pt]
Check if variable empty
  & \texttt{if "\%VAR\%"==""}
  & \texttt{if [ -z "\$VAR" ]} \\[4pt]
Check if directory exists
  & \texttt{if exist "DIR\textbackslash"}
  & \texttt{if [ -d "DIR" ]} \\[4pt]
Exit a script/function
  & \texttt{goto :eof}
  & \texttt{return} (in function/sourced script) \\
\bottomrule
\end{tabular}
\end{table}

\begin{table}[h]
\centering
\caption{\texttt{uv} quick reference: equivalences with \texttt{pip} and \texttt{venv}}
\label{tab:uv-cheatsheet}
\begin{tabular}{@{}p{0.32\textwidth}p{0.28\textwidth}p{0.28\textwidth}@{}}
\toprule
\textbf{Task} & \textbf{pip / venv} & \textbf{uv} \\
\midrule
Install \texttt{uv}
  & (not applicable)
  & See Section~\ref{sec:uv} \\[4pt]
Create environment
  & \texttt{python -m venv PATH}
  & \texttt{uv venv PATH -{}-python 3.12} \\[4pt]
Activate environment
  & \texttt{source PATH/bin/activate}
  & \texttt{source PATH/bin/activate} \\[4pt]
Install a package
  & \texttt{pip install pkg}
  & \texttt{uv pip install pkg} \\[4pt]
Install from \texttt{pyproject.toml}
  & \texttt{pip install -e .}
  & \texttt{uv sync} \\[4pt]
Add a dependency
  & Edit \texttt{requirements.txt} manually
  & \texttt{uv add pkg} \\[4pt]
Remove a dependency
  & Edit \texttt{requirements.txt} manually
  & \texttt{uv remove pkg} \\[4pt]
Generate lockfile
  & \texttt{pip freeze > req.txt}
  & \texttt{uv lock} (or implicit in \texttt{uv sync}) \\[4pt]
Install from lockfile
  & \texttt{pip install -r req.txt}
  & \texttt{uv sync} \\[4pt]
Launcher script (Windows)
  & \texttt{ActivateEnv.bat}
  & \texttt{ActivateUVEnv.bat} \\[4pt]
Launcher function (Linux)
  & \texttt{ActivateEnv()} in \texttt{.bashrc}
  & \texttt{ActivateUVEnv()} in \texttt{.bashrc} \\[4pt]
Environment root
  & \texttt{C:\textbackslash penv} / \texttt{\textasciitilde/penv}
  & \texttt{C:\textbackslash uvenv} / \texttt{\textasciitilde/uvenv} \\
\bottomrule
\end{tabular}
\end{table}

% ============================================================
\section{Summary}
\label{sec:summary}
% ============================================================

This chapter covered the tools and techniques for managing Python
environments and automating project setup across Windows and Linux:

\begin{itemize}
  \item \textbf{Python virtual environments}: isolated, per-project
        copies of Python that prevent dependency conflicts between
        projects. Each environment has its own set of installed packages,
        and deleting one has no effect on any other.
  \item \textbf{The standard \texttt{pip} and \texttt{venv} toolchain}:
        the foundation-backed tools included with every Python
        installation for creating environments, installing packages, and
        recording dependencies in \texttt{requirements.txt}. Knowing
        this toolchain remains an essential skill regardless of what
        other tools you adopt.
  \item \textbf{\texttt{uv}}: a modern, high-performance alternative
        developed by Astral that offers dramatically faster installs, a
        global package cache, reproducible lockfiles with cryptographic
        hashes, and native \texttt{pyproject.toml} support. Because
        \texttt{uv} automatically synchronises dependencies from a
        lockfile via \texttt{uv sync}, environment reproduction becomes
        a single command rather than a manual process.
  \item \textbf{Level~1 scripting}: writing reusable launcher scripts
        (\texttt{ActivateEnv.bat} and \texttt{ActivateEnv.sh} for
        \texttt{pip}-managed environments; \texttt{ActivateUVEnv.bat}
        and the \texttt{ActivateUVEnv()} shell function for
        \texttt{uv}-managed environments) and project-specific
        one-liners (\texttt{\ProjectName{}.bat}, \texttt{\ProjectName{}.sh}) that
        activate the correct environment, change to the project
        directory, and launch VS Code in a single invocation.
  \item \textbf{Level~2 scripting}: shell functions defined in
        \texttt{\textasciitilde/.bashrc} that provide a one-word launch
        command from any directory, with no \texttt{source} prefix
        required.
  \item \textbf{The \texttt{source} vs.\ execute distinction}: on
        Linux, running a script normally spawns a child shell whose
        environment changes vanish when it exits. The \texttt{source}
        command (or its shorthand \texttt{.}) runs the script in the
        current shell, which is why virtual environment activation
        requires sourcing. This same distinction explains why
        \texttt{ActivateEnv.sh} uses \texttt{return} instead of
        \texttt{exit} for error handling.
\end{itemize}

The \texttt{\ProjectName{}} command you built in this chapter is a small but
meaningful example of a broader principle: the CLI lets you encode your
workflow into repeatable, shareable, and improvable scripts. Every
tedious setup step you eliminate today is cognitive load freed up for the
work that actually matters.
